\documentclass[10pt,twocolumn,letterpaper]{article}

\usepackage{times}
\usepackage[margin=1in,columnsep=0.25in]{geometry}
\usepackage{amsmath}
\usepackage{amssymb}
\usepackage{graphicx}
\usepackage{url}
\usepackage{hyperref}
\usepackage{abstract}
\usepackage{titlesec}

\hypersetup{colorlinks=true, urlcolor=black, linkcolor=black, citecolor=black}

\titleformat{\section}{\normalfont\bfseries}{\thesection.}{0.5em}{}
\titleformat{\subsection}{\normalfont\itshape}{\thesubsection}{0.5em}{}

\renewcommand{\abstractnamefont}{\normalfont\bfseries}
\renewcommand{\abstracttextfont}{\normalfont\small\itshape}
\setlength{\absleftindent}{0pt}
\setlength{\absrightindent}{0pt}

\title{\Large\bf Cookie Provenance: A Compression-Based Approach\\to Affiliate Fraud Detection}

\author{
  Satoshi Nakamoto \\
  \texttt{satoshin@gmx.com} \\
  \texttt{www.bitcoin.org}
}

\date{}

\begin{document}

\twocolumn[
  \maketitle
  \begin{abstract}
    Affiliate cookie stuffing is commission fraud. A site places tracking
    cookies in a browser without a click by the user, then collects a
    commission when the user makes a purchase, and the merchant pays for
    a referral that never happened. We propose a detection system based on
    a principle from lossless data compression: a domain that sets a cookie
    but has never appeared in a user's navigation history is a novelty event
    in the sense of Lempel and Ziv, and a high rate of such events is evidence
    of stuffing. We combine this novelty signal with five supporting signals
    derived from request metadata to produce a suspicion score for each domain,
    and we implement the full system as a browser extension that runs on the
    user's device and transmits nothing externally. Results on two commercial
    sites show novelty rates of 88 percent and fraud scores of 0.80 for
    confirmed affiliate networks. We characterize the conditions under which
    detection degrades and the cost an adversary must pay to evade it.
  \end{abstract}
  \vspace{1em}
]

\section{Introduction}

Commerce on the web depends on affiliate marketing. A publisher earns a
commission when a user follows a referral link and completes a purchase, and
the commission is evidence that the publisher drove the sale. Cookie stuffing
corrupts this. A site places an affiliate cookie without any click, and when
the user later completes a purchase, the stuffing party collects a commission
they did not earn.

The scale of the problem is established by recent events. In 2026, Bloomberg
reported that 51 percent of credited merchandise value at a major shopping
application came from cookie stuffing~\cite{bloomberg2026}. Impact.com
suspended the company. Its operators said they did not know.

We build a tool to make that claim harder to sustain. The core idea is
borrowed from Lempel and Ziv, who showed in 1978 that a sequence is
compressible to the extent that it can be encoded as references to substrings
seen before~\cite{lz78}. We apply this to a browsing session. The set of
domains a user has navigated to is a dictionary. A cookie from a domain
outside that dictionary is a novelty miss. A high miss rate means something
is placing cookies that the user's navigation did not cause.

\section{The Problem of Cookie Provenance}

We define $\mathcal{N}$ to be the set of hostnames a user has navigated to
over their browsing history, and we call $\mathcal{N}$ the navigation
dictionary. We define $\mathcal{E} = \{e_1, e_2, \ldots, e_n\}$ to be the
sequence of cookie-setting events in a session. Each event $e_i$ carries a
domain, a cookie name, a request URL, a resource type, a timestamp in
milliseconds from page load, and a Referer header if one was sent.

A cookie-setting event is legitimate if it follows causally from a user's
act of navigation. It is illegitimate if it does not. We cannot observe
the causal relationship directly, so we approximate it with signals described
in Section~\ref{sec:signals}.

We define the LZ novelty indicator for domain $d$ as
\[
  \lambda(d) =
  \begin{cases}
    0 & \text{if } d \in \mathcal{N} \text{ or } \mathrm{apex}(d) \in \mathcal{N}, \\
    1 & \text{otherwise,}
  \end{cases}
\]
where $\mathrm{apex}(d)$ is the registrable domain of $d$. This is the
primary signal. Five signals support it.

\section{Signal Architecture}
\label{sec:signals}

We assign each domain $d$ a suspicion score $s(d) \in [0, 1]$:
\begin{align}
  s(d) = \min\bigl(1,\;
    &0.30 \cdot \lambda(d) \nonumber \\
    &+ 0.25 \cdot \alpha_{\mathrm{url}}(d) \nonumber \\
    &+ 0.15 \cdot \alpha_{\mathrm{ck}}(d) \nonumber \\
    &+ 0.15 \cdot \rho(d) \nonumber \\
    &+ 0.10 \cdot \tau(d) \nonumber \\
    &+ 0.05 \cdot \nu(d) \bigr).
\end{align}
Each signal takes a value in $[0, 1]$. We aggregate events from the same
domain by taking the maximum observed value of each signal, for reasons
discussed in Section~\ref{sec:aggregation}.

\textbf{LZ novelty} $\lambda(d)$, weight 0.30. A domain absent from
navigation history has no ordinary explanation for its cookies. This is
the dominant signal and the hardest to defeat.

\textbf{Affiliate URL pattern} $\alpha_{\mathrm{url}}(d)$, weight 0.25.
Score is 1 if the request URL matches known affiliate tracking patterns,
including path segments \texttt{/click}, \texttt{/track},
\texttt{/redirect}, and query parameters \texttt{affiliate\_id},
\texttt{aff\_id}, \texttt{ref}, \texttt{subid}, \texttt{clickid}.
Confirmed network domains receive this score unconditionally.

\textbf{Affiliate cookie name} $\alpha_{\mathrm{ck}}(d)$, weight 0.15.
Score is 1 if the cookie name matches patterns common to affiliate systems,
including prefixes \texttt{aff}, \texttt{affiliate}, \texttt{partner},
\texttt{ref}, and network-specific prefixes \texttt{cj\_}, \texttt{sa\_},
\texttt{impact\_}.

\textbf{Hidden resource type} $\rho(d)$, weight 0.15. Score is 1.0 for
image, XHR, beacon, ping, and websocket resources; 0.7 for script; 0.0
for navigation. These types carry tracking pixels and invisible cookie
requests that users cannot see or initiate.

\textbf{Early timing} $\tau(d)$, weight 0.10. A cookie set within 500
milliseconds of page load cannot follow a click. Score decays as
\[
  \tau(d) =
  \begin{cases}
    1.0 & t < 0.5\text{s,} \\
    \max\!\left(0,\, 1 - \frac{t - 0.5}{1.5}\right) & 0.5 \leq t < 2.0\text{s,} \\
    0.0 & t \geq 2.0\text{s.}
  \end{cases}
\]

\textbf{No Referer} $\nu(d)$, weight 0.05. This is the fraction of cookies
from $d$ arriving without a Referer header. Legitimate referrals carry one.
Absence removes a link in provenance.

\subsection{Weight Justification}
\label{sec:weights}

Weights reflect the specificity of each signal to stuffing versus ordinary
browsing. LZ novelty is the most discriminating: a domain absent from
navigation history has no ordinary explanation for cookies it sets. URL and
cookie name patterns are highly specific but depend on the operator using
known patterns. The remaining signals are supporting evidence. Weight
calibration against a labeled dataset is in progress; the values above
reflect an informed prior.

\section{The Navigation Dictionary}

Each time the browser commits to a URL in a top-level frame, we add the
hostname and its apex domain to $\mathcal{N}$. The dictionary persists
across sessions in browser local storage and grows monotonically.

The dictionary is retroactive. When a user navigates to domain $d$,
$\lambda(d)$ is set to 0 for all existing reports and scores are
recomputed. This prevents false positives from domains that set cookies
before a page finishes loading.

As $|\mathcal{N}|$ grows, the prior probability of a legitimate domain
being absent decreases. A novelty miss from a mature dictionary is stronger
evidence than one from a fresh dictionary.

\section{The Affiliate Network Fingerprint}

We maintain a database of 17 affiliate networks mapped to their tracking
domains: Commission Junction, ShareASale, Awin, Rakuten, Impact,
PartnerStack, Partnerize, ClickBank, FlexOffers, PepperJam, Tradedoubler,
Viglink, Skimlinks, MaxBounty, Refersion, Tune, HasOffers, and Google
DoubleClick. A domain matching this database is classified as a source of
fraud. A domain that scores above threshold without matching the database is
classified as suspicious. The distinction matters for remediation and for
regulatory framing.

\section{Signal Aggregation}
\label{sec:aggregation}

We aggregate signals at the domain level by taking the maximum observed
value across all cookie events from that domain. Stuffing manifests in at
least one high-signal event; taking the mean would let a stuffing operator
dilute detection by mixing legitimate and illegitimate requests. The maximum
is conservative in the sense that it is least likely to miss a detection.

The cost is sensitivity to single anomalous requests. In practice, the
combination of independent signals limits this risk. A domain that scores
high on LZ novelty alone, without matching URL or cookie name patterns,
is more likely a cold-start artifact than a fraud source.

\section{Empirical Results}

We ran the detector on two sites in August 2026. Navigation history
contained approximately 3,000 domains at the time of testing.

On Slickdeals.net, 53 domains set cookies in a single session. Of these,
47 had no entry in the navigation dictionary. The novelty rate was 88
percent. These 47 domains set cookies through embedded resources despite
the user having no prior relationship with them.

On Yahoo.com, \texttt{doubleclick.net}, operated by Google and a member
of the fingerprint database, received a score of 0.80. The cookie arrived
via an image resource within 400 milliseconds of page load, with no Referer
header. All four non-novelty signals contributed. The domain was confirmed
as fraud by the fingerprint database.

Both results are reproducible. Any Firefox user who installs the extension
and visits these sites will observe comparable behavior.

\section{Adversarial Analysis}

A stuffing operator who knows the detection architecture has several options.

The most complete defense against LZ novelty is to age a domain into users'
navigation histories before using it to stuff. This requires genuine traffic,
which is costly. It also does not defeat the other five signals.

The timing signal can be defeated by delaying cookie-setting past two
seconds. This constrains stuffing to pages with high dwell time. Cookie
name and URL signals can be defeated with opaque identifiers, but this
eliminates two signals without affecting the other four. The Referer signal
can be defeated by injecting a synthetic header server-side, but this is not
observable from the browser.

LZ novelty is the hardest signal to defeat. To suppress it, an operator
must cause a user to genuinely navigate to the cookie-setting domain. The
cost of defeating the primary signal is the cost of becoming a site people
visit.

\section{Privacy}

The system runs entirely on the user's device. Navigation history and domain
reports live in browser local storage and are never transmitted. The
extension observes network traffic to score events but does not log URLs,
cookie values, or page content. The only action visible outside the browser
is deletion of cookies at user request, via \texttt{browser.cookies.remove()}.

Private browsing sessions are excluded from the navigation dictionary.

\section{Conclusion}

We have described a detection system for affiliate cookie stuffing, built on
LZ novelty scoring and five supporting signals derived from request metadata.
It runs locally, transmits nothing, and requires no account or external
resource.

A company valued at 180 million dollars credited 51 percent of its sales to
cookies that were stuffed. Its operators said they did not know. We have built
the instrument that tests this claim, and we have made it available to any
user of Firefox at no cost.

The browser does not lie. The signal is there for anyone who looks.

\begin{thebibliography}{9}

\bibitem{lz78}
  A.~Lempel and J.~Ziv,
  ``Compression of Individual Sequences via Variable-Rate Coding,''
  \textit{IEEE Transactions on Information Theory}, vol.~24, no.~5,
  pp.~530--536, 1978.

\bibitem{nakamoto}
  S.~Nakamoto,
  ``Bitcoin: A Peer-to-Peer Electronic Cash System,''
  2008. [Online]. Available: \url{https://bitcoin.org/bitcoin.pdf}

\bibitem{bloomberg2026}
  Bloomberg News,
  ``Phia Cookie Stuffing Investigation,''
  July 2026.

\bibitem{cookiestuff}
  quantumcelnav,
  ``CookieStuff: Real-time Affiliate Cookie Stuffing Detection,''
  2026. [Online]. Available: \url{https://github.com/quantumcelnav/cookiestuff}

\end{thebibliography}

\end{document}
